% ============================================================
\setpartpreamble{\partintro
  {Wie LINQ zu SQL wird, wie ein Plugin eigene Funktionen und Member übersetzt und wie aus
  Zeilen wieder Objekte werden.}
  {
    \item Kapitel~\ref{ch:the-query-pipeline}: Die Query-Pipeline
    \item Kapitel~\ref{ch:translating-methods-and-members}: Methoden und Member übersetzen
    \item Kapitel~\ref{ch:preventing-client-side-evaluation}: Clientseitige Auswertung
    verhindern
    \item Kapitel~\ref{ch:further-query-extension-points}: Weitere Erweiterungspunkte der
    Abfrage
    \item Kapitel~\ref{ch:materialization-from-rows-to-objects}: Materialisierung: von Zeilen zu Objekten
  }}
\part{Abfragen}\label{part:IV}

% ------------------------------------------------------------
\chapter{Die Query-Pipeline}\label{ch:the-query-pipeline}

Eine LINQ-Abfrage durchläuft in EF Core mehrere Stufen, bevor SQL entsteht. Wer weiß, wo
welche Stufe sitzt, weiß auch, wo ein Plugin ansetzen muss:

\begin{flowdiagram}
  \flowstep{Funcletizer}{Wertet Teilausdrücke ohne Datenbezug in .NET aus und macht
  Parameter daraus. Plugin: IEvaluatableExpressionFilterPlugin}
  \flowstep{Preprocessing}{Navigationen expandieren, Query-Filter einfügen, normalisieren.
  Plugin: IQueryExpressionInterceptor}
  \flowstep{Queryable-Übersetzung}{Where, Select, OrderBy usw. werden zu einer
  SelectExpression}
  \flowstep{SQL-Übersetzung}{Ausdrücke werden zu SqlExpressions. Plugins: Method-, Member-
  und Aggregate-Translators}
  \flowstep{Postprocessing}{Nullability-Optimierung, Parameter-abhängige Anpassungen}
  \flowstep{SQL-Generierung}{QuerySqlGenerator schreibt den SQL-Text}
  \flowstep{Materialisierung}{Zeilen werden zu Objekten (Kapitel~\ref{ch:materialization-from-rows-to-objects}). Plugin:
  IMaterializationInterceptor}
\end{flowdiagram}

Die kompilierte Abfrage wird gecacht. Translators laufen also nur einmal pro Abfrageform,
nicht bei jeder Ausführung. Werte, die sich pro Ausführung ändern (Parameter,
Session-Settings), müssen deshalb im SQL als Parameter oder als Datenbankfunktion stehen,
nie als zur Übersetzungszeit eingesetzte Konstante. Genau deshalb übersetzt das Beispiel
\code{Current\allowbreak{}Tenant(\allowbreak{})} in einen Aufruf von
\code{current\_\allowbreak{}setting(\allowbreak{}.\allowbreak{}.\allowbreak{}.\allowbreak{})}
und setzt nicht den aktuellen Mandanten ein.

% ------------------------------------------------------------
\chapter{Methoden und Member übersetzen}\label{ch:translating-methods-and-members}

\section{Das EF.Functions-Muster}

Datenbankfunktionen, die keine Entsprechung in .NET haben, werden als Extension-Methods auf
\code{Db\allowbreak{}Functions} deklariert. Sie werden nie ausgeführt; der Körper wirft nur
eine Exception für den Fall, dass jemand sie außerhalb von LINQ aufruft. Die Deklaration
liegt im Kernpaket, die Übersetzung im Provider-Paket:

\begin{sourcecode}[Tenancy\allowbreak{}Db\allowbreak{}Functions\allowbreak{}Extensions.\allowbreak{}cs (\allowbreak{}Kernpaket)]{csharp}
namespace Microsoft.EntityFrameworkCore;

public static class TenancyDbFunctionsExtensions
{
    public static TenantId CurrentTenant(this DbFunctions _)
        => throw new InvalidOperationException(
            $"{nameof(CurrentTenant)} ist nur in LINQ-to-Entities-Abfragen verwendbar.");
}
\end{sourcecode}

\section{Ein Method-Call-Translator}

Ein \code{IMethod\allowbreak{}Call\allowbreak{}Translator} bekommt jeden Methodenaufruf, den
EF nicht selbst kennt. Er prüft, ob er zuständig ist, und baut dann aus
\code{ISql\allowbreak{}Expression\allowbreak{}Factory} einen SQL-Ausdrucksbaum:

\begin{sourcecode}[Query/\allowbreak{}Current\allowbreak{}Tenant\allowbreak{}Translator.\allowbreak{}cs]{csharp}
public class CurrentTenantTranslator : IMethodCallTranslator
{
    private static readonly MethodInfo CurrentTenantMethod =
        typeof(TenancyDbFunctionsExtensions).GetRuntimeMethod(
            nameof(TenancyDbFunctionsExtensions.CurrentTenant), [typeof(DbFunctions)])!;

    private readonly ISqlExpressionFactory _sql;
    private readonly IRelationalTypeMappingSource _typeMappingSource;
    private readonly ITenancySingletonOptions _options;

    public CurrentTenantTranslator(
        ISqlExpressionFactory sql,
        IRelationalTypeMappingSource typeMappingSource,
        ITenancySingletonOptions options)
    {
        // ... ArgumentNullException.ThrowIfNull für alle Parameter ...
        _sql = sql;
        _typeMappingSource = typeMappingSource;
        _options = options;
    }

    public virtual SqlExpression? Translate(
        SqlExpression? instance, MethodInfo method,
        IReadOnlyList<SqlExpression> arguments,
        IDiagnosticsLogger<DbLoggerCategory.Query> logger)
    {
        if (method != CurrentTenantMethod)
            return null;                          // nicht zuständig -> nächster Translator

        var text = _typeMappingSource.FindMapping(typeof(string))!;
        var boolean = _typeMappingSource.FindMapping(typeof(bool))!;

        var currentSetting = _sql.Function(
            "current_setting",
            [
                _sql.Constant(_options.Settings.SettingName, text),
                _sql.Constant(true, boolean),
            ],
            nullable: true,
            argumentsPropagateNullability: [true, false],
            typeof(string), text);

        var nullIfEmpty = _sql.Function(
            "nullif",
            [currentSetting, _sql.Constant(string.Empty, text)],
            nullable: true,
            argumentsPropagateNullability: [true, false],
            typeof(string), text);

        return _sql.Convert(nullIfEmpty, typeof(TenantId), TenantIdTypeMapping.Default);
    }
}
\end{sourcecode}

Aus \code{o.\allowbreak{}Tenant\allowbreak{}Id ==
EF.\allowbreak{}Functions.\allowbreak{}Current\allowbreak{}Tenant(\allowbreak{})} wird
damit:

\begin{sourcecode}[Erzeugtes SQL (\allowbreak{}Ausschnitt)]{sql}
WHERE o."TenantId" = CAST(nullif(current_setting('app.tenant_id', TRUE), '') AS uuid)
\end{sourcecode}

\begin{datatable}{H{0.60}H{1.40}}
  \textbf{Detail} & \textbf{Bedeutung} \\
  \tablehead
  \code{return null} & Vertrag für „nicht zuständig“. Niemals eine Exception für fremde
  Methoden werfen; dann ist der nächste Translator an der Reihe. \\
  \tablerule
  \code{arguments} & Bei Extension-Methods ist \code{arguments[0]} der \code{this}-Parameter
  (hier \code{Db\allowbreak{}Functions}). \\
  \tablerule
  \code{nullable} & Kann das Ergebnis NULL sein? EF optimiert damit \code{IS
  NULL}-Prüfungen. Falsche Angaben führen zu falschen Ergebnissen. \\
  \tablerule
  \code{arguments\allowbreak{}Propagate\allowbreak{}Nullability} & Wird das Ergebnis NULL,
  wenn dieses Argument NULL ist? Für \code{nullif(\allowbreak{}x,\allowbreak{} '')}: ja für
  x, nein für den Leerstring. \\
  \tablerule
  Type Mappings & Jeder SQL-Ausdruck braucht ein Mapping. Es bestimmt Literale, Parameter
  und Casts. \code{Convert} mit dem TenantId-Mapping erzeugt
  \code{CAST(\allowbreak{}.\allowbreak{}.\allowbreak{}.\allowbreak{} AS uuid)}. \\
  \tablerule
  \code{Constant(\allowbreak{}.\allowbreak{}.\allowbreak{}.\allowbreak{})} & Konstanten
  werden als Literal über das Mapping erzeugt und korrekt maskiert, auch der Setting-Name.
  \\
\end{datatable}

\section{Ein Member-Translator}

Für Properties und Felder gibt es \code{IMember\allowbreak{}Translator}. Das Beispiel
übersetzt \code{tenant\allowbreak{}Id.\allowbreak{}Is\allowbreak{}Empty}. Die Konstante
bekommt das Type Mapping der Spalte, damit das Literal im richtigen Format entsteht:

\begin{sourcecode}[Query/\allowbreak{}Tenant\allowbreak{}Id\allowbreak{}Member\allowbreak{}Translator.\allowbreak{}cs]{csharp}
public class TenantIdMemberTranslator : IMemberTranslator
{
    private static readonly PropertyInfo IsEmptyProperty =
        typeof(TenantId).GetProperty(nameof(TenantId.IsEmpty))!;

    private readonly ISqlExpressionFactory _sql;

    public TenantIdMemberTranslator(ISqlExpressionFactory sql)
    {
        ArgumentNullException.ThrowIfNull(sql);

        _sql = sql;
    }

    public virtual SqlExpression? Translate(
        SqlExpression? instance, MemberInfo member, Type returnType,
        IDiagnosticsLogger<DbLoggerCategory.Query> logger)
    {
        if (instance is null || !Equals(member, IsEmptyProperty))
            return null;

        return _sql.Equal(instance, _sql.Constant(TenantId.Empty, instance.TypeMapping));
    }
}
\end{sourcecode}

\section{Translator-Plugins}

Translators werden nicht direkt registriert, sondern über Plugin-Klassen, die eine Liste
liefern. Die Plugins sind \textbf{Scoped} (Anhang~\ref{app:A}) und können daher
Scoped-Dienste wie \code{ISql\allowbreak{}Expression\allowbreak{}Factory} injizieren:

\begin{sourcecode}[Query/\allowbreak{}Tenancy\allowbreak{}Translator\allowbreak{}Plugins.\allowbreak{}cs]{csharp}
public class TenancyMethodCallTranslatorPlugin : IMethodCallTranslatorPlugin
{
    public TenancyMethodCallTranslatorPlugin(
        ISqlExpressionFactory sql,
        IRelationalTypeMappingSource typeMappingSource,
        ITenancySingletonOptions options)
    {
        Translators = [new CurrentTenantTranslator(sql, typeMappingSource, options)];
    }

    public virtual IEnumerable<IMethodCallTranslator> Translators { get; }
}

public class TenancyMemberTranslatorPlugin : IMemberTranslatorPlugin
{
    public TenancyMemberTranslatorPlugin(ISqlExpressionFactory sql)
    {
        Translators = [new TenantIdMemberTranslator(sql)];
    }

    public virtual IEnumerable<IMemberTranslator> Translators { get; }
}
\end{sourcecode}

\begin{tip}
  \code{query.\allowbreak{}To\allowbreak{}Query\allowbreak{}String(\allowbreak{})} erzeugt
  das SQL einer Abfrage, ohne sie auszuführen und ohne Datenbankverbindung. Es ist das
  wichtigste Werkzeug, um Translators zu entwickeln und zu testen
  (Kapitel~\ref{ch:test-strategy}).
\end{tip}

% ------------------------------------------------------------
\chapter{Clientseitige Auswertung verhindern}\label{ch:preventing-client-side-evaluation}

Der erste Schritt der Pipeline, der Funcletizer, wertet jeden Teilausdruck vorab in .NET
aus, der keinen Bezug zu Abfragedaten hat. Aus \code{o.\allowbreak{}Total > limit * 2} wird
so ein Parameter. Für
\code{EF.\allowbreak{}Functions.\allowbreak{}Current\allowbreak{}Tenant(\allowbreak{})} wäre
das fatal: Der Ausdruck hat keinen Datenbezug, also würde EF die Methode ausführen, und die
wirft ihre Exception.

EF kennt nur wenige Ausnahmen. Im Quellcode von
\code{Evaluatable\allowbreak{}Expression\allowbreak{}Filter} (EF Core 10) sind das
\code{Date\allowbreak{}Time.\allowbreak{}Now},
\code{Date\allowbreak{}Time.\allowbreak{}Utc\allowbreak{}Now},
\code{Date\allowbreak{}Time.\allowbreak{}Today}, die
\code{Date\allowbreak{}Time\allowbreak{}Offset}-Gegenstücke,
\code{Guid.\allowbreak{}New\allowbreak{}Guid(\allowbreak{})},
\code{Random.\allowbreak{}Next(\allowbreak{}.\allowbreak{}.\allowbreak{}.\allowbreak{})} und
alle Methoden, die in EFs eigener Klasse
\code{Db\allowbreak{}Functions\allowbreak{}Extensions} deklariert sind. Eigene
EF.Functions-Methoden liegen in einer anderen Klasse und fallen \textbf{nicht} darunter.
Danach fragt EF alle registrierten Plugins:

\begin{sourcecode}[Kernpaket: Tenancy\allowbreak{}Db\allowbreak{}Functions\allowbreak{}Attribute.\allowbreak{}cs,\allowbreak{} Query/\allowbreak{}Tenancy\allowbreak{}Evaluatable…Plugin.\allowbreak{}cs]{csharp}
[AttributeUsage(AttributeTargets.Class, Inherited = false)]
public sealed class TenancyDbFunctionsAttribute : Attribute;

public class TenancyEvaluatableExpressionFilterPlugin : IEvaluatableExpressionFilterPlugin
{
    public virtual bool IsEvaluatableExpression(Expression expression)
        => expression is not MethodCallExpression { Method.DeclaringType: { } type }
           || !type.IsDefined(typeof(TenancyDbFunctionsAttribute), inherit: false);
}

[TenancyDbFunctions]
public static class TenancyDbFunctionsExtensions { /* CurrentTenant() */ }
\end{sourcecode}

Das Plugin liegt im \textbf{Kernpaket} und erkennt Funktionsklassen an einem Attribut. So
deckt ein einziges Plugin auch die PostgreSQL-spezifischen Funktionen aus
Kapitel~\ref{ch:custom-sql-expression-nodes} und \ref{ch:aggregate-translators} und jedes
künftige Provider-Paket ab, ohne dass der Kern diese Klassen kennen muss.

\begin{important}[Wann das Plugin nötig ist]
  Ob eine eigene EF.Functions-Methode dieses Plugin braucht, hängt davon ab, ob sie
  Argumente mit Datenbezug hat.
  \code{EF.\allowbreak{}Functions.\allowbreak{}Soundex(\allowbreak{}c.\allowbreak{}Name)}
  ist unproblematisch, weil \code{c.\allowbreak{}Name} auf eine Spalte zeigt. Methoden ohne
  Argumente oder nur mit Konstanten (\code{Current\allowbreak{}Tenant(\allowbreak{})},
  \code{Random(\allowbreak{})}, \code{Server\allowbreak{}Time(\allowbreak{})}) brauchen es
  immer. Fehlt es, kommt die Exception aus dem Methodenkörper, was schwer zu diagnostizieren
  ist.
\end{important}

% ------------------------------------------------------------
\chapter{Weitere Erweiterungspunkte der Abfrage}\label{ch:further-query-extension-points}

Das Beispiel braucht Method- und Member-Translators. Komplexe Plugins greifen auf weitere
Punkte zurück; hier ein Überblick mit Einordnung:

\begin{datatable}{H{0.96}H{1.08}H{0.96}}
  \textbf{Erweiterungspunkt} & \textbf{Zweck} & \textbf{Hinweise} \\
  \tablehead
  \code{IAggregate\allowbreak{}Method\allowbreak{}Call\allowbreak{}Translator\allowbreak{}Plugin}
  & Eigene Aggregatfunktionen in \code{Group\allowbreak{}By} (z. B.
  \code{bool\_\allowbreak{}and}, Statistikfunktionen) & Scoped, mehrfach registrierbar;
  umgesetzt in Kapitel~\ref{ch:aggregate-translators} \\
  \tablerule
  \code{IQuery\allowbreak{}Expression\allowbreak{}Interceptor} & LINQ-Ausdrucksbaum vor der
  Übersetzung umschreiben & Singleton-Interceptor; ideal für Query-Rewrites, Tags, Hints \\
  \tablerule
  \code{IMaterialization\allowbreak{}Interceptor} & Objekte bei der Materialisierung
  erzeugen, initialisieren oder prüfen & Singleton-Interceptor; muss zustandslos und
  thread-sicher sein; umgesetzt in Kapitel~\ref{ch:materialization-from-rows-to-objects} \\
  \tablerule
  \code{IMethod\allowbreak{}Call\allowbreak{}Translator\allowbreak{}Provider} u. a. &
  Translator-Infrastruktur ersetzen & Einzeldienste des Providers; nur in Ausnahmefällen \\
  \tablerule
  \code{IQuery\allowbreak{}Sql\allowbreak{}Generator\allowbreak{}Factory} & SQL-Generierung
  für eigene SqlExpression-Typen & Singleton, Provider-Dienst; umgesetzt in
  Kapitel~\ref{ch:custom-sql-expression-nodes} \\
  \tablerule
  \code{IRelational\allowbreak{}Parameter\allowbreak{}Based\allowbreak{}Sql\allowbreak{}Processor\allowbreak{}Factory}
  & Nullability-Verarbeitung eigener Ausdrucksknoten & Scoped, Provider-Dienst; umgesetzt in
  Kapitel~\ref{ch:custom-sql-expression-nodes} \\
  \tablerule
  \code{IQuery\allowbreak{}Translation\allowbreak{}Postprocessor\allowbreak{}Factory} &
  Ausdrucksbaum nach der Übersetzung anpassen & Provider-Dienst; tiefer Eingriff, eng an
  EF-Interna gebunden \\
\end{datatable}

\begin{note}[Vorkompilierte Abfragen]
  Mit NativeAOT und vorkompilierten Abfragen (seit EF Core 9) übersetzt EF Abfragen bereits
  beim Build. Translators laufen dann im Build-Prozess, nicht zur Laufzeit. Sie dürfen
  deshalb nur von Options und Modell abhängen, nie von Laufzeitzustand. Das Beispiel erfüllt
  das: Der Setting-Name ist eine Option, der Mandant selbst wird erst in der
  Datenbank-Session ausgewertet.
\end{note}

% ------------------------------------------------------------
\chapter{Materialisierung: von Zeilen zu Objekten}\label{ch:materialization-from-rows-to-objects}

Die bisherigen Kapitel dieses Teils beschreiben den Hinweg: aus LINQ wird SQL.
Kapitel~\ref{ch:the-update-pipeline} beschreibt den Weg von Entitäten zu INSERT, UPDATE und
DELETE. Dieses Kapitel schließt den Kreis und beschreibt den Rückweg: Wie werden aus den
Zeilen, die die Datenbank liefert, Objekte, und wo kann ein Plugin dabei eingreifen?

\section{Der Rückweg im Überblick}

\begin{flowdiagram}
  \flowstep{SQL-Ausführung}{Der Befehl läuft; ein Command-Interceptor sieht ihn vorher und
  nachher. Plugin: IDbCommandInterceptor}
  \flowstep{DbDataReader}{Eine Zeile nach der anderen, Spalten in der Reihenfolge der
  Projektion. EF-intern}
  \flowstep{Shaping}{Der kompilierte Shaper liest die Spalten, wandelt Werte über das Type
  Mapping um und baut das Ergebnis. EF-intern, gesteuert über Modell und Type Mappings}
  \flowstep{Identity-Prüfung}{Beim Tracking: Ist die Entität schon im Change Tracker, gilt
  die vorhandene Instanz. EF-intern}
  \flowstep{Materializer}{Konstruktor aufrufen, Properties und Felder setzen. Plugin:
  IMaterializationInterceptor; Modell: Constructor Binding}
  \flowstep{CLR-Instanz}{Das fertige Objekt}
  \flowstep{Tracking}{Die Instanz wird Unchanged im Change Tracker; Ereignis
  ChangeTracker.Tracked. Anwendung}
  \flowstep{Relationship Fixup}{Navigationen zu bereits geladenen Entitäten werden
  verbunden. EF-intern}
  \flowstep{Abfrageergebnis}{Liste, Enumerator oder Einzelwert}
\end{flowdiagram}

Die Identity-Prüfung steht \emph{vor} dem Materializer: Eine Entität, die der Context
bereits trackt, wird nicht ein zweites Mal erzeugt. Wo ein Plugin eingreifen kann und was
EF für sich behält, zeigt die folgende Übersicht:

\begin{datatable}{H{0.62}H{1.24}H{1.14}}
  \textbf{Stufe} & \textbf{Eingriff für Plugins} & \textbf{EF-intern} \\
  \tablehead
  Ausführung & \code{IDb\allowbreak{}Command\allowbreak{}Interceptor} kann den Befehl ändern
  oder einen eigenen \code{Db\allowbreak{}Data\allowbreak{}Reader} liefern
  (Kapitel~\ref{ch:interceptors}) & Verbindung, Ausführungsstrategie, Wiederholungen \\
  \tablerule
  Shaping & Indirekt: Type Mapping und Value Converter bestimmen, wie Werte gelesen und
  umgewandelt werden (Kapitel~\ref{ch:type-mappings-and-value-converters}) & Aufbau und
  Kompilierung des Shapers \\
  \tablerule
  Materializer & \code{IMaterialization\allowbreak{}Interceptor} mit vier Haken;
  Constructor Binding über das Modell & Wie EF Konstruktor und Werte zusammensetzt \\
  \tablerule
  Tracking & Ereignisse \code{Tracked} und \code{State\allowbreak{}Changed} des Change
  Trackers (Anwendung) & Identity Map, Zustände, Schnappschüsse für die Änderungserkennung
  \\
  \tablerule
  Fixup & Kein Haken & Navigation Fixer \\
\end{datatable}

\section{Shaping: der kompilierte Rückweg}

Beim Übersetzen einer Abfrage entsteht nicht nur SQL, sondern auch ein \textbf{Shaper}: ein
Ausdrucksbaum, der beschreibt, wie aus einer Zeile des
\code{Db\allowbreak{}Data\allowbreak{}Reader} das Ergebnis wird. EF kompiliert ihn zu einem
Delegaten und cacht ihn zusammen mit dem SQL pro Abfrageform
(Kapitel~\ref{ch:the-query-pipeline}). Pro Zeile läuft nur noch dieser Delegat, ohne
Nachschlagen im Modell.

Das hat eine wichtige Folge: Alles, was den Shaper prägt, wird beim Übersetzen entschieden,
nicht beim Ausführen. Dazu gehören die Lesemethode des Type Mappings, der Value Converter,
der gewählte Konstruktor und die registrierten Materialization-Interceptors. Für
\code{Tenant\allowbreak{}Id} liest der Shaper den Providerwert (\code{Guid}) und wendet den
Ausdruck des \code{Tenant\allowbreak{}Id\allowbreak{}Converter} an, der in den Shaper
einkompiliert ist. Werte in JSON-Spalten liest er über den
\code{Json\allowbreak{}Value\allowbreak{}Reader\allowbreak{}Writer} des Mappings. Kein
Plugin-Code muss dafür pro Wert laufen; deshalb ist ein Value Converter für reine
Wertumwandlung immer billiger als ein Interceptor.

Eine Projektion wie \code{Select(\allowbreak{}o => o.\allowbreak{}Total)} erzeugt einen
Shaper, der nur diese Spalte liest. Es entsteht keine Entität, und kein Materializer läuft.

\section{Der Materializer: Instanz, Konstruktor, Felder}

Für eine Entität erzeugt der Shaper zuerst eine Instanz und setzt dann die übrigen Werte.
Welchen Konstruktor er verwendet, legt das Modell fest (Constructor Binding):

\begin{datatable}{H{0.65}H{1.35}}
  \textbf{Regel} & \textbf{Verhalten in EF Core 10} \\
  \tablehead
  Parameterloser Konstruktor & Ist einer vorhanden, auch ein privater, verwendet EF ihn. \\
  \tablerule
  Konstruktor mit Parametern & Sonst einer, dessen Parameter nach Name und Typ zu gemappten
  Properties passen. Diese Properties setzt EF danach nicht noch einmal. \\
  \tablerule
  Dienste als Parameter & \code{Db\allowbreak{}Context}, \code{ILazy\allowbreak{}Loader} und
  \code{IEntity\allowbreak{}Type} kann EF übergeben. \\
  \tablerule
  Übrige Properties & EF schreibt sie über das Backing-Field, wenn es eines gibt. Logik im
  Setter läuft beim Laden deshalb \emph{nicht}. \\
  \tablerule
  Shadow-Properties & Haben kein Feld; beim Tracking liegen ihre Werte im Change Tracker. \\
  \tablerule
  Owned Types & Sind Entity-Typen und werden einzeln materialisiert. \\
  \tablerule
  Complex Types & Gehören zur Entität und werden mit ihr materialisiert, ohne eigenen
  Durchlauf. \\
\end{datatable}

Dass Setter beim Laden nicht laufen, ist gewollt: Laden ist kein Erzeugen. Prüfungen, die
beim Anlegen oder Ändern greifen sollen, gehören in Konstruktor und Methoden der Domäne. Was
beim Laden gilt, garantiert die Datenbank.

\section{Tracking, Identity Resolution und Fixup}

Was nach dem Materializer passiert, hängt vom Tracking-Verhalten der Abfrage ab:

\begin{datatable}{H{0.80}H{1.10}H{1.10}}
  \textbf{Modus} & \textbf{Identität} & \textbf{Fixup} \\
  \tablehead
  Tracking (Standard) & Vor dem Materializer prüft EF den Schlüssel im Change Tracker. Bei
  einem Treffer gibt es die vorhandene Instanz zurück, unverändert: Lokale Änderungen
  bleiben, neue Werte aus der Datenbank werden verworfen. & Neue Instanzen werden Unchanged;
  EF verbindet Navigationen zu allen getrackten Entitäten, auch aus früheren Abfragen. \\
  \tablerule
  \code{As\allowbreak{}No\allowbreak{}Tracking} & Jede Zeile ergibt eine neue Instanz, auch
  dieselbe Entität mehrfach. & Nur innerhalb des Ergebnisses: Der Shaper verbindet, was über
  \code{Include} geladen wurde. \\
  \tablerule
  \code{As\allowbreak{}No\allowbreak{}Tracking\allowbreak{}With\allowbreak{}Identity\allowbreak{}Resolution}
  & Eine Identity Map nur für diese Abfrage: Dieselbe Entität wird einmal materialisiert. &
  Wie ohne Tracking; danach trackt niemand die Instanzen. \\
\end{datatable}

\textbf{Relationship Fixup} ist EF-intern. Ein Plugin kann es nicht beeinflussen, eine
Anwendung kann es über \code{Change\allowbreak{}Tracker.\allowbreak{}Tracked} und
\code{State\allowbreak{}Changed} beobachten.

\begin{important}[IIdentityResolutionInterceptor ist kein Materialisierungs-Haken]
  Der Name legt es nahe, aber \code{IIdentity\allowbreak{}Resolution\allowbreak{}Interceptor}
  (Scoped, mehrfach) läuft bei Abfragen nie. Er greift, wenn
  \code{Attach}, \code{Update} oder \code{Add} eine Instanz mit einem Schlüssel anhängen
  wollen, den der Context schon trackt, und entscheidet, was mit der getrackten Instanz
  geschieht. EF liefert dafür
  \code{Ignoring\allowbreak{}Identity\allowbreak{}Resolution\allowbreak{}Interceptor} und
  \code{Updating\allowbreak{}Identity\allowbreak{}Resolution\allowbreak{}Interceptor} mit.
  Abfragen lösen Identität ohne ihn: beim Tracking über den Change Tracker, sonst wie in der
  Tabelle oben.
\end{important}

\section{IMaterializationInterceptor im Detail}

\code{IMaterialization\allowbreak{}Interceptor} hat vier Haken, die EF für jede Entität in
dieser Reihenfolge aufruft. Alle haben eine Standardimplementierung; ein Plugin überschreibt
nur, was es braucht.

\begin{datatable}{H{0.75}H{1.05}H{1.20}}
  \textbf{Methode} & \textbf{Zeitpunkt} & \textbf{Möglichkeiten} \\
  \tablehead
  \code{Creating\allowbreak{}Instance} & Vor dem Konstruktor & Werte der Zeile lesen, auch
  Shadow-Properties; mit \code{Suppress\allowbreak{}With\allowbreak{}Result} eine eigene
  Instanz liefern \\
  \tablerule
  \code{Created\allowbreak{}Instance} & Nach dem Konstruktor & Über den Konstruktor gesetzte
  Werte sind da; eine andere Instanz zurückgeben \\
  \tablerule
  \code{Initializing\allowbreak{}Instance} & Vor dem Setzen der übrigen Werte & Mit
  \code{Suppress} das Setzen selbst übernehmen \\
  \tablerule
  \code{Initialized\allowbreak{}Instance} & Alle Werte gesetzt & Prüfen, ergänzen oder eine
  andere Instanz zurückgeben \\
\end{datatable}

Jede Methode bekommt
\code{Materialization\allowbreak{}Interception\allowbreak{}Data}: den Context, den
\code{IEntity\allowbreak{}Type}, das Tracking-Verhalten der Abfrage und
\code{Get\allowbreak{}Property\allowbreak{}Value} für die Werte der Zeile. Weil der Typ ein
\code{IEntity\allowbreak{}Type} ist, laufen Complex Types nicht einzeln durch den
Interceptor, Owned Types dagegen schon. Der Konstruktor der Struktur ist EF-intern; Tests
rufen deshalb nicht die Haken auf, sondern die Logik dahinter (siehe unten).

\begin{important}[Lebensdauer: ein Singleton, einkompiliert in den Shaper]
  Der Interceptor ist ein \code{ISingleton\allowbreak{}Interceptor}. EF baut ihn beim
  Übersetzen in den Shaper ein, und der Shaper wird pro internem Service Provider gecacht.
  Eine Instanz bedient deshalb alle Contexts und alle Threads gleichzeitig; sie muss
  zustandslos und threadsicher sein. Daten einer Anfrage holt sie bei jedem Aufruf über
  \code{materialization\allowbreak{}Data.\allowbreak{}Context}.

  Singleton-Interceptor-Instanzen gehören außerdem zum Cache-Schlüssel des internen
  Providers. Wer pro Context eine neue Instanz über
  \code{Add\allowbreak{}Interceptors} übergibt, bekommt pro Context einen neuen internen
  Provider samt Modell; ab dem 21. bricht EF standardmäßig mit
  \code{Many\allowbreak{}Service\allowbreak{}Providers\allowbreak{}Created\allowbreak{}Warning}
  ab. Ein Plugin registriert ihn in \code{Apply\allowbreak{}Services}; eine Anwendung
  übergibt stets dieselbe Instanz.
\end{important}

\section{Beispiel: Mandanten beim Laden prüfen}

Kapitel~\ref{ch:interceptors} zeigt eine echte Lücke: Übergibt die Anwendung eine bereits
geöffnete Verbindung, läuft der Connection-Interceptor nicht, und eine Abfrage liefert still
die Zeilen des Mandanten, dessen Wert noch in der Session steht. Dasselbe passiert, wenn die
Anwendung versehentlich mit einer Rolle verbindet, die RLS umgeht. Das Plugin bietet dafür
eine optionale Prüfung an, die genau an der richtigen Stelle sitzt: Jede mandantenfähige
Entität, die eine Abfrage materialisiert, muss zum Mandanten des Contexts gehören.

\begin{sourcecode}[Anwendung]{csharp}
options.UseNpgsql(connectionString)
       .UseAcmeTenancy(sp, tenancy => tenancy.VerifyMaterializedTenants());
\end{sourcecode}

\begin{sourcecode}[Interceptors/\allowbreak{}Tenant\allowbreak{}Materialization\allowbreak{}Interceptor.\allowbreak{}cs (\allowbreak{}Kernpaket,\allowbreak{} gekürzt)]{csharp}
public class TenantMaterializationInterceptor : IMaterializationInterceptor
{
    public virtual InterceptionResult<object> CreatingInstance(
        MaterializationInterceptionData materializationData,
        InterceptionResult<object> result)
    {
        Verify(materializationData.Context, materializationData.EntityType,
            materializationData.GetPropertyValue);

        return result;
    }

    internal static void Verify(
        DbContext context, IEntityType entityType, Func<string, object?> readValue)
    {
        // Laufzeit-Annotation der Isolations-Convention; ohne sie keine Isolation
        if (entityType.GetTenantPropertyName() is not { } propertyName)
            return;

        // Bei jeder Zeile nachschlagen, nie merken: Diese Instanz teilen sich alle Contexts
        var current = FindTenantProvider(context).CurrentTenantId;
        var value = (TenantId)readValue(propertyName)!;

        if (value != current)
            throw new TenancyException("... Row level security did not take effect ...");
    }

    // FindTenantProvider: über die Interceptors des Contexts (wie ApplyTenantSession) zum
    // ITenantProvider, den der SaveChanges-Interceptor dieses Contexts trägt
}
\end{sourcecode}

\begin{sourcecode}[Infrastructure/\allowbreak{}Tenancy\allowbreak{}Options\allowbreak{}Extension.\allowbreak{}cs › Apply\allowbreak{}Services (\allowbreak{}Ergänzung)]{csharp}
// Optionale Prüfung beim Laden: ein Singleton-Interceptor im internen Provider, eine
// Instanz pro Provider. Die Einstellung steht im Settings-Record und damit im Hash.
if (Settings.VerifyMaterializedTenants)
    new EntityFrameworkServicesBuilder(services)
        .TryAdd<ISingletonInterceptor, TenantMaterializationInterceptor>();
\end{sourcecode}

\begin{datatable}{H{0.70}H{1.30}}
  \textbf{Entscheidung} & \textbf{Begründung} \\
  \tablehead
  \code{Creating\allowbreak{}Instance} statt
  \code{Initialized\allowbreak{}Instance} & Die Prüfung liest den Wert aus der Zeile. Eine
  fremde Zeile wird nie zum Objekt, und die Mandanten-Property darf eine Shadow-Property
  sein. \\
  \tablerule
  Laufzeit-Annotation & Welche Entität isoliert ist und wie ihre Mandanten-Property heißt,
  steht seit dem Modellbau am Entity-Typ
  (Kapitel~\ref{ch:property-level-fluent-api-and}). Kein Raten zur Laufzeit. \\
  \tablerule
  Mandant pro Aufruf & Über den Context zum \code{ITenant\allowbreak{}Provider} dieses
  Contexts. Kostet pro Entität ein Nachschlagen und einen Vergleich; deshalb ist die Prüfung
  optional. \\
  \tablerule
  Kein Mandant & Ohne Mandanten liefert RLS keine Zeilen. Kommt trotzdem eine, wirkt RLS
  nicht, und die Prüfung schlägt an. \\
  \tablerule
  Im Kernpaket & Die Prüfung kennt weder SQL noch Provider-Typen
  (Kapitel~\ref{ch:what-belongs-in-the-core}); beide Provider-Pakete aktivieren sie mit einer
  Zeile. \\
\end{datatable}

Die Integrationstests stellen die Lücke aus Kapitel~\ref{ch:interceptors} nach: Ein Context
für Mandant A über eine Verbindung, deren Session noch Mandant B trägt. Ohne Prüfung liefert
die Abfrage die Zeilen von B, mit Prüfung wirft sie eine
\code{Tenancy\allowbreak{}Exception}. Eigene Zeilen passieren mit und ohne Tracking und
über \code{From\allowbreak{}Sql}.

\section{Der Fehlerfall: Anfragedaten im Singleton}

Der naheliegende erste Entwurf übergibt dem Interceptor den Mandanten im Konstruktor und
erzeugt ihn wie die anderen Interceptors pro Context. Das scheitert am Cache-Schlüssel: Jeder
Context bekommt einen eigenen internen Provider. Der naheliegende zweite Entwurf teilt eine
Instanz und merkt sich den Mandanten beim ersten Aufruf, um das Nachschlagen pro Zeile zu
sparen:

\begin{sourcecode}[tests/\allowbreak{}.\allowbreak{}.\allowbreak{}.\allowbreak{}Integration\allowbreak{}Tests/\allowbreak{}Materialization\allowbreak{}Tests.\allowbreak{}cs (\allowbreak{}FALSCH,\allowbreak{} gekürzt)]{csharp}
private sealed class RememberingTenantInterceptor : IMaterializationInterceptor
{
    private TenantId? _tenant;                       // Anfragedaten in einem Singleton

    public InterceptionResult<object> CreatingInstance(
        MaterializationInterceptionData materializationData,
        InterceptionResult<object> result)
    {
        // "Optimierung": den Mandanten nur einmal ermitteln
        _tenant ??= _tenantOf(materializationData.Context);

        if (materializationData.GetPropertyValue<TenantId>(nameof(Order.TenantId)) != _tenant)
            throw new TenancyException("Row of another tenant.");

        return result;
    }
}
\end{sourcecode}

Der Test zeigt beide Folgen. Die zweite Anfrage für Mandant B bekommt einen Fehler für ihre
eigenen Zeilen, weil der Interceptor sich Mandant A gemerkt hat. Schlimmer: Ein Context für
B über eine Verbindung, deren Session noch A trägt, liefert die Zeilen von A, und der
gemerkte Mandant A lässt sie passieren. Die Prüfung verhindert das Leck nicht, sie verdeckt
es. Unter Last kommt eine Race Condition hinzu, denn alle Threads schreiben dasselbe Feld.

\begin{tip}
  Ein Singleton-Interceptor hat nur zwei Quellen für Daten: was zur Konfiguration gehört
  (über Konstruktor oder Singleton-Options, für alle Contexts gleich) und was der Aufruf
  mitbringt (\code{Context}, Entity-Typ, Werte der Zeile). Alles andere gehört in die
  Anwendung oder in einen Scoped-Dienst, den der Interceptor über den Context erreicht.
\end{tip}

\section{Grenzen}

\begin{datatable}{H{0.65}H{1.35}}
  \textbf{Grenze} & \textbf{Folge} \\
  \tablehead
  Nur Entitäten & Projektionen einzelner Spalten oder in DTOs laufen nicht durch den
  Materializer. Ein Integrationstest belegt, dass
  \code{Select(\allowbreak{}o => o.\allowbreak{}Total)} auch Werte eines fremden Mandanten
  ungeprüft liefert; schützen kann hier nur die Datenbank. \\
  \tablerule
  Keine getrackten Entitäten & Ist eine Entität schon getrackt, wird sie nicht erneut
  materialisiert; der Interceptor sieht nur die erste Abfrage. \\
  \tablerule
  Kein Filter & Zeilen auslassen kann der Interceptor nicht, nur werfen oder die Instanz
  ersetzen. Filtern ist Aufgabe von Query-Filter oder Datenbank. \\
  \tablerule
  Kein SQL & Er sieht das Ergebnis, nicht die Abfrage; \code{Execute\allowbreak{}Update},
  \code{Execute\allowbreak{}Delete} und Skalarabfragen materialisieren nichts. \\
  \tablerule
  Heißer Pfad & Er läuft für jede Entität jeder Abfrage
  (Kapitel~\ref{ch:plugin-performance}); teure Arbeit vervielfacht sich mit der
  Ergebnisgröße. \\
  \tablerule
  Exception bricht die Abfrage ab & Auch alle übrigen Zeilen sind verloren. Das ist für eine
  Sicherheitsprüfung richtig, für eine Bequemlichkeitsfunktion nicht. \\
\end{datatable}

\section{Wann Materialisierung, wann etwas anderes?}

Der Interceptor ist der richtige Ort für Querschnittsaufgaben, die \emph{jede} geladene
Entität betreffen, über alle Abfragen, mit und ohne Tracking, auch bei rohem SQL. Für fast
alles andere gibt es einen einfacheren und günstigeren Weg:

\begin{datatable}{H{1.00}H{0.95}H{1.05}}
  \textbf{Ziel} & \textbf{Besser} & \textbf{Warum} \\
  \tablehead
  Einen Wert anders speichern als darstellen (Format, Typ, Verschlüsselung) & Value
  Converter & Wirkt in beide Richtungen und in Parametern, einkompiliert statt pro Aufruf \\
  \tablerule
  Unveränderliche Objekte, Invarianten & Constructor Binding & EF ruft den Konstruktor
  selbst auf \\
  \tablerule
  Nur einige Spalten oder ein DTO & Projektion mit \code{Select} & Keine Entität, kein
  Tracking, weniger Daten \\
  \tablerule
  Dieselbe Einstellung für viele Typen (Konverter, Feldzugriff) & Convention & Einmal beim
  Modellbau statt bei jeder Zeile \\
  \tablerule
  Abgeleitete Werte, Geschäftsregeln & Domänenlogik (Methode, berechnete Property) & Laden
  ist kein Geschäftsvorfall \\
  \tablerule
  Reagieren, wenn eine Anwendung Entitäten lädt & \code{Change\allowbreak{}Tracker.\allowbreak{}Tracked}
  & Öffentliche API ohne Plugin; nur bei Tracking \\
  \tablerule
  Prüfung oder Ergänzung jeder geladenen Entität, auch ohne Tracking &
  \code{IMaterialization\allowbreak{}Interceptor} & Der einzige Haken, der alle Wege zur
  Entität sieht \\
\end{datatable}
